Geographic gaming --- manipulating the resolution level or adjacency graph definition
to produce a partisan outcome while maintaining algorithmic neutrality as cover ---
is the most subtle attack vector against algorithmic redistricting.
We test two geographic gaming scenarios: resolution choice gaming (varying the
$k/n$ threshold that determines tract vs.\ block-group resolution) and adjacency
manipulation (selectively removing cross-partisan tract boundaries from the
adjacency graph).

For resolution gaming: varying the threshold from 4\% to 6\% (a plausible
dispute range) changes the resolution choice for 3--4 states and produces a
maximum national partisan shift of 0.18 Democratic seats.

For adjacency manipulation: a synthetic ``poisoned'' NC adjacency graph that
excludes 50 strategically chosen cross-partisan tract edges (the maximum
manipulation feasible without violating contiguity) produces a partisan shift
of 0.15 Democratic seats.

Both gaming vectors are bounded well below strategic significance ($< 0.5$ seats)
and both are detectable via the \texttt{bisect} audit chain:
resolution choice is recorded in the plan manifest (\texttt{plan\_resolution}),
and adjacency modifications produce a different SHA-256 hash from the
standard TIGER/Line adjacency graph.
The audit chain converts the geographic gaming threat from a covert manipulation
to a detectable deviation.
